
\begin{abstract}

Software maps are a commonly used tool for code quality monitoring in software-development projects and decision making processes. 
While providing an important visualization technique for the hierarchical system structure of a single software revision, they lack capabilities with respect to the visualization of changes over multiple revisions. 
This paper presents a novel technique for visualizing the evolution of the software system structure based on software metric trends. 
These trend maps extend software maps by using real-time rendering techniques for natural phenomena yielding additional visual variables that can be effectively used for the communication of changes. 
Therefore, trend data is automatically computed by hierarchically aggregating software metrics. 
We demonstrate and discuss the presented technique using two real world data sets of complex software systems.

\begin{classification} % according to http://www.acm.org/class/1998/
	\CCScat{Software Engineering}{D.2.2}{Design Tools and Techniques}{Computer-aided software engineering (CASE)}
	\CCScat{Software Engineering}{D.2.7}{Distribution, Maintenance, and Enhancement}{Documentation}
	\CCScat{Computer Graphics}{I.3.7}{Three-Dimensional Graphics and Realism }{Color, shading, shadowing, and texture}
\end{classification}

\end{abstract}
